\documentstyle[%


righttag%


,11pt%


,amssymb%


,russian%               remove in Eng.


,izvru%                 Set izven% in Eng.


% ,izvru-ks             for brief communications


]{amsart}





%       Math definitions





\newcommand{\tts}[1]{}





\theoremstyle{definition}


\theorembodyfont{\rm}


\newtheorem{notenonum}{\hskip\parindent Замечание}


\renewcommand{\thenotenonum}{}








%\hoffset=-15mm                  For Eng.





\setcounter{page}{24}


\god{1997}


\nomer{10~(425)} %                Remove in Eng.


%\nomer{Vol.\,41, No.\,10}        Set in Eng.


% \str{75-81}                    Set in Eng.


\udk{517.518}


\author{\it В.Э.\,Гейт}


\title{Об аппроксимационных свойствах\\


высших производных периодических функций}





\date{}


% \pagestyle{myheadings}         Set in Eng.





\pagestyle{plain} %              Remove in Eng.


\begin{document}





% \thispagestyle{plain}          Set in Eng.





\maketitle





% Body text





Каждая монотонно убывающая к нулю последовательность $\{F_n\}^\infty_{n=0}$


задает класс $F_C$


непрерывных на $(-\infty,+\infty)$ периода $2\pi$ функций $f$, наилучшие


тригонометрические


приближения которых удовлетворяют условию $E_n(f)\le F_n$


($n=0,1,\dotsc$). Будем


говорить, что порядок величины $\alpha_n$ есть $\beta_n$ и писать


$\alpha_n\sim\beta_n$, если $\alpha_n=O(\beta_n)$ и


$\beta_n=O(\alpha_n)$.





\begin{thm}


Для произвольно заданного $r$ $(=1,2,\dotsc)$ и любой последовательности


$F_n\downarrow 0$ $(n\uparrow\infty)$ по функциям $f\in F_C$ имеем


$$


\sup E_n(f^{(r)})\sim n^rF_n+\sum^\infty_{\nu=n+1}\nu^{r-1}F_\nu


$$


в предположении сходимости указанного ряда.


\end{thm}





{\bf Доказательство.} Пусть $f\in F_C$ и ряд


$\sum\nu^{r-1}E_\nu(f)<+\infty$ сходится,


тогда


(напр., \cite{Bar}, c.\,488) $r$-я производная $f^{(r)}$ непрерывна,


причем


\begin{equation}


E_n(f^{(r)})=O\Bigl(n^rE_n(f)+\sum^\infty_{\nu=n+1}\nu^{r-1}


E_\nu(f)\Bigr).


\end{equation}


Отсюда по функциям $f$ класса $F_C$ имеем с очевидностью


\begin{equation}


\sup E_n(f^{(r)})=O\Bigl(n^rF_n+\sum^\infty_{\nu=n+1}


\nu^{r-1}F_\nu \Bigr).


\end{equation}





Обратное к (2) $O$-соотношение докажем следующим образом. По данному


$r$


($=1,2,\dotsc$) и последовательности $F_n\downarrow 0$, полагая


$\Delta F_{\nu-1}=F_{\nu-1}-F_\nu$ ($\nu=1,2,\dotsc$),


рассмотрим функцию вида


\begin{equation}


f(x)=f_{r,F}(x)=\sum^\infty_{\nu=1}\Delta F_{\nu-1}\cos\Bigl(


\nu x+r\frac{\pi}{2}\Bigr).


\end{equation}


Так как уклонение функции (3) от ее частичной суммы $S_n(f,x)$ есть


$$


f(x)-S_n(f,x)=\sum^\infty_{\nu=n+1}\Delta F_{\nu-1}\cos\Bigl(


\nu x+r\frac{\pi}{2}\Bigr),


$$


то из-за предположения $F_\nu\downarrow 0$ будем иметь


\begin{equation}


E_n(f)\le\sum^\infty_{\nu=n+1}\Delta F_{\nu-1}=F_n\quad


(n=0,1,\dotsc).


\end{equation}


В силу сходимости ряда (3) к $f\in F_C$ нетрудно видеть, что $r$-я


производная функции (3) непрерывна, причем


\begin{equation}


f^{(r)}(x)=(-1)^r\sum^\infty_{\nu=1}\nu^r\Delta F_{\nu-1}


\cos\nu x.


\end{equation}


Поскольку коэффициенты из (5) неотрицательны, то по теореме Ченя


\cite{Cen} будем иметь


$$


\|f^{(r)}-S_n(f^{(r)})\|_C=O(E_n(f^{(r)})).


$$


Отсюда ввиду (5)


\begin{equation}


\sum^\infty_{\nu=n+1}\nu^r\Delta F_{\nu-1}=O(E_n(f^{(r)})).


\end{equation}


Но для левой части в (6) имеем последовательно


\begin{gather*}


\sum^\infty_{\nu=n+1}\nu^r\Delta F_{\nu-1}>


\sum^\infty_{\nu=n+1}\Delta F_{\nu-1}\biggl(\,\sum^n_{j=1}j^{r-1}+


\sum^\nu_{j=n+1}j^{r-1}\biggr)=


F_n\sum^n_{j=1}j^{r-1}+\sum^\infty_{\nu=n+1}\Delta F_{\nu-1}


\sum^\nu_{j=n+1}j^{r-1}> \\


%


> r^{-1}F_nn^r+\sum^\infty_{j=n+1}j^{r-1}F_{j-1}\ge r^{-1}


\biggl(n^rF_n+\sum^\infty_{j=n+1}j^{r-1}F_j\biggr).


\end{gather*}


Последнее неравенство совместно c (6), а также с (4) доказывает, что


\begin{equation}


n^rF_n+\sum^\infty_{\nu=n+1}\nu^{r-1}F_\nu=


O(\sup E_n(f^{(r)}))


\end{equation}


по функциям $f\in F_C$. Теорема 1 ввиду (2) и (7) доказана.





Ниже через $\Tilde f$ обозначается тригонометрически сопряженная с $f$


функция.





\begin{thm}


Пусть задано $r$ $(=1,2,\dotsc)$ и последовательность $F_n\downarrow 0$


при $n\uparrow\infty$, тогда по функциям $f$ класса $F_C$ имеем


\begin{equation}


\sup E_n(\Tilde f^{(r)})\sim n^rF_n+\sum^\infty_{\nu=n+1}


\nu^{r-1}F_\nu


\end{equation}


в предположении сходимости ряда в {\em (8)}.


\end{thm}





Доказательство проводится аналогично доказательству теоремы 1.





Через $\omega_k(f,\delta)$ ($0<\delta\le\pi$) будем далее обозначать


$k$-й модуль непрерывности функции $f$


(см., напр., \cite{Bar}, c.\,487).





\begin{thm}


Для произвольно заданных натуральных $r$, $k$ и последовательности


$F_n\downarrow 0$ $(n\uparrow\infty)$


по функциям $f\in F_C$ имеем


$$


\sup\omega_k(f^{(r)},n^{-1})\sim\biggl(n^{-k}


\sum^k_{\nu=1}\nu^{k+r-1}F_{\nu-1}+\sum^\infty_{\nu=n+1}


\nu^{r-1}F_\nu\biggr)


$$


в предположении сходимости указанного здесь ряда.


\end{thm}





\begin{pf} Ввиду соотношения (см., напр, \cite{Bar}, c.\,490)


\begin{equation}


\omega_k(f^{(r)},n^{-1})=O\biggl(n^{-k}\sum^n_{\nu=1}


\nu^{k+r-1}E_{\nu-1}(f)+\sum^\infty_{\nu=n+1}\nu^{r-1}E_\nu(f)


\biggr)


\end{equation}


по функциям $f\in F_C$ имеем


\begin{equation}


\sup\omega_k(f^{(r)},n^{-1})=O\biggl(n^{-k}\sum^n_{\nu=1}


\nu^{k+r-1}F_{\nu-1}+\sum^\infty_{\nu=n+1}\nu^{r-1}F_\nu\biggr).


\end{equation}


Из (7) и обобщенного неравенства Джексона имеем


\begin{equation}


n^rF_n+\sum^\infty_{\nu=n+1}\nu^{r-1}F_\nu=O(\sup\omega_k


(f^{(r)},n^{-1})).


\end{equation}


После этого необходимо также показать, что первое слагаемое в (10) есть


$O$ от $\sup\omega_k(f^{(r)},n^{-1})$ по классу $F_C$. Для этого по


заданным


$k$, $r$ ($=1,2,\dotsc$) и по данной последовательности


$F_n\downarrow 0$


($n\uparrow\infty$) рассмотрим функцию


\begin{equation}


f(x)=f_{r,F,k}(x)=\sum^\infty_{\nu=1}\Delta F_{\nu-1}


\cos\Bigl(\nu x-(k+r)\frac{\pi}{2}\Bigr),


\end{equation}


очевидно лежащую в классе $F_C$. Тогда из-за сходимости ряда


$\sum\limits^\infty_{\nu=1}\nu^{r-1}F_\nu$ $r$-я


производная функции (12) непрерывна (\cite{Bar}, c.\,488), причем легко


видеть, что


\begin{equation}


f^{(r)}(x)=(-1)^r\sum^\infty_{\nu=1}\nu^r\Delta F_{\nu-1}


\cos\Bigl(\nu x-k\frac{\pi}{2}\Bigr).


\end{equation}


Так что при $k$ нечетном (13) --- это синус-ряд, а при $k$ четном ---


косинус-ряд. В таком случае для оценки снизу $k$-го модуля непрерывности


функции (13) применимо неравенство (23) леммы 3 работы автора (\cite{Ge1},


c.\,72), по которому


\begin{equation}


n^{-k}\sum^n_{\nu=1}\nu^{k+r}\Delta F_{\nu-1}=


O(\omega_k(f^{(r)},n^{-1})).


\end{equation}


Так как $\nu^{k+r}\ge\sum\limits^\nu_{j=1}j^{k+r-1}$, то сумма из левой


части (14) превосходит величину


$$


\sum^n_{\nu=1}\Delta F_{\nu-1}\sum^\nu_{j=1}j^{k+r-1}


=\sum^n_{j=1}j^{k+r-1}\sum^n_{\nu=j}\Delta F_{\nu-1}=


\sum^n_{j=1}j^{k+r-1}(F_{j-1}-F_n)>\sum^n_{j=1}j^{k+r-1}


F_{j-1}-F_n n^{k+r}.


$$


Поэтому для функции (12) имеем сначала первое $O$-соотношение


$$


n^{-k}\sum^n_{j=1}j^{k+r-1}F_{j-1}=O(\omega_k(f^{(r)},n^{-1}))


+n^rF_n=O\Bigl(\sup_{f\in F_C}\omega_k(f^{(r)},n^{-1})\Bigr),


$$


а второе $O$-соотношение получается с учетом (11).


\end{pf}





Аналогично получается





\begin{thm}


Пусть заданы натуральные $k$, $r$ и последовательность


$F_n\downarrow 0$ $(n\uparrow \infty)$, тогда по


функциям $f$ класса $F_C$ имеем


\begin{equation}


\sup\omega_k(\Tilde f^{(r)},n^{-1})\sim n^{-k}\sum^n_{\nu=1}


\nu^{k+r-1}F_{\nu-1}+\sum^\infty_{\nu=n+1}\nu^{r-1}F_\nu


\end{equation}


при условии сходимости ряда в {\em (15)}.


\end{thm}





Далее, имея функцию $\omega(\delta)>0$ ($0<\delta\le\pi$) с условием


$\varliminf\omega(\delta)=0$


при $\delta\to+0$, рассмотрим класс $H_l(\omega)$,


состоящий из тех непрерывных на $(-\infty,+\infty)$ периода $2\pi$


функций $f$, у которых $\omega_l(f,\delta)\le\omega(\delta)$


($0<\delta\le\pi$). Исправленную мажоранту порядка $l$ функции


$\omega(\delta)$


обозначим через $\omega^{**}_l(\delta)$ (см. \cite{St}). Имеет место





\begin{thm}


Пусть заданы натуральные $l$, $r$ и мажоранта $\omega(\delta)$, при


которых сходится ряд


\begin{equation}


\sum^\infty_{\nu=1}\nu^{r-1}\omega^{**}_l(\nu^{-1})<+\infty,


\end{equation}


тогда при $\omega^{**}_l(\delta)\not\equiv 0$ необходимо $l>r$, причем


\begin{equation}


\sup E_n(f^{(r)})\sim \sum^\infty_{\nu=n+1}\nu^{r-1}


\omega^{**}_l(\nu^{-1}),\quad f\in H_l(\omega),


\end{equation}


кроме случая\,\footnote{Это ограничение со слов ``кроме случая$\dotsc$'',


как и в теоремах 6, 7, 8, устраняется леммой 1 \cite{Ge3}. Аналогично


можно опустить последнюю строку теорем 9, 10.}, когда оба $l$, $r$ четные.


\end{thm}


\addtocounter{footnote}{-1}





{\bf Доказательство.} Исключив сразу тривиальный случай


$\omega^{**}_l(\delta)\equiv 0$ на $(0,\pi]$,


согласно \cite{St} будем иметь


\begin{equation}


0<\omega^{**}_l(\delta)\downarrow 0 \quad (\delta\downarrow +0);\qquad


\delta^{-l}\omega^{**}_l(\delta)\downarrow \quad (\delta\uparrow).


\end{equation}


В частности, последовательность $\nu^l\omega^{**}_l(\nu^{-1})$ возрастает


при $\nu\uparrow$, откуда


очевидным является факт расходимости ряда (16) при $r\ge l$. Так как


согласно


(\cite{Bar}, c.\,488)


$$


E_n(f^{(r)})=O\biggl(n^rE_n(f)+\sum^\infty_{\nu=n+1}\nu^{r-1}


E_\nu(f)\biggr),


$$


то по обобщенному неравенству Джексона


$E_\nu(f)=O(\omega_l(f,\nu^{-1}))$


имеем


$$


E_n(f^{(r)})=O\biggl(n^r\omega_l(f,n^{-1})+\sum^\infty_{\nu=n+1}


\nu^{r-1}\omega_l(f,\nu^{-1})\biggr).


$$


В силу общих свойств модулей непрерывности (\cite{Bar}, c.\,487) нетрудно


вывести, что первое слагаемое здесь есть $O$ от второго, так что


окончательно


\begin{equation}


E_n(f^{(r)})=O\biggl(\,\sum^\infty_{\nu=n+1}\nu^{r-1}\omega_l


(f,\nu^{-1})\biggr)


\end{equation}


в предположении сходимости ряда в (19). Пусть теперь


$f\in H_l(\omega)$, тогда, в


частности,


$\omega_l(f,\nu^{-1})\le\omega(\nu^{-1})$ ($\nu=1,2,\dotsc$) и согласно


\cite{St} выполнено также соотношение


$\omega_l(f,\nu^{-1})=O(\omega^{**}_l(\nu^{-1}))$. Поэтому в силу (19)


\begin{equation}


\sup E_n(f^{(r)})=O\biggl(\,\sum^\infty_{\nu=n+1}\nu^{r-1}


\omega^{**}_l(\nu^{-1})\biggr),\quad f\in H_l(\omega).


\end{equation}


Получим теперь $O$-соотношение, обратное к (20), кроме случая, когда


$l$,


$r$ оба четные. Пусть $r$ нечетное, а $l$ произвольное с условием


$l>r \ge 1$.


Рассмотрим функцию $f_3$ из работы автора \cite{Ge1} (см. там (38) при


$k=l$ и $\varphi_\nu=\omega^{**}_l(\nu^{-1})$), т.\,е.


\begin{equation}


f_3(x)=\sum^\infty_{\nu=1}a^{(l)}_\nu\sin\nu x,\quad


a^{(l)}_\nu=\varphi_\nu-\Bigl(1-\frac{1}{\nu}\Bigr)^l\varphi_{\nu-1}.


\end{equation}


В силу возрастания последовательности $\nu^l\varphi_\nu$ все


$a^{(l)}_\nu\ge 0$ и $\omega_l(f_3,n^{-1})=O(\omega^{**}_l(n^{-1}))$


согласно


(\cite{Ge1}, сс.\,75, 76). Отсюда, как известно (\cite{Bar}, c.\,487),


$\omega_l(f,\delta)=O(\omega^{**}_l(\delta))$


($0<\delta\le\pi$). Поэтому при некотором $c>0$ функция


$cf_3\in H_l(\omega)$. Вместе


с тем $r$-я производная функции (21) при нечетном числе $r$ есть


\begin{equation}


f^{(r)}_3(x)=\pm\sum^\infty_{\nu=1}\nu^ra^{(l)}_\nu\cos\nu x.


\end{equation}


Значит, согласно \cite{Cen} для наилучшего приближения функции (22) ввиду


$a^{(l)}_\nu\ge 0$ имеем


\begin{equation}


\|f^{(r)}_3-S_n(f^{(r)}_3)\|=O(E_n(f^{(r)}_3)),\quad


\sum^\infty_{\nu=n+1}\nu^ra^{(l)}_\nu=O(E_n(f^{(r)}_3)).


\end{equation}


Согласно (\cite{Ge1}, с.\,75, (40)) коэффициенты ряда (21) допускают


представление


$$


a^{(l)}_\nu=-\Delta\varphi_{\nu-1}+l\varphi_{\nu-1}/\nu


+O(\varphi_{\nu-1}/\nu^2).


$$


Поэтому левая часть (23) имеет вид


\begin{equation}


-\sum^\infty_{\nu=n+1}\nu^r\Delta\varphi_{\nu-1}+l


\sum^\infty_{\nu=n+1}\nu^{r-1}\varphi_{\nu-1}+O\biggl(\,


\sum^\infty_{\nu=n+1}\nu^{r-2}\varphi_{\nu-1}\biggr).


\end{equation}


Для третьего слагаемого здесь очевидным является неравенство


\begin{equation}


\sum^\infty_{\nu=n+1}\nu^{r-2}\varphi_\nu\le\frac{1}{n+1}


\sum^\infty_{\nu=n+1}\nu^{r-1}\varphi_{\nu-1}=o\biggl(\,


\sum^\infty_{\nu=n+1}\nu^{r-1}\varphi_{\nu-1}\biggr).


\end{equation}


Далее, т.\,к. $\nu^r\le r\sum\limits^\nu_{j=1}j^{r-1}$, то сумма из


первых слагаемых в (24) допускает оценку сверху величиной


\begin{multline*}


r\sum^\infty_{\nu=n+1}\Delta\varphi_{\nu-1}\biggl(\,


\sum^n_{j=1}j^{r-1}+\sum^\nu_{j=n+1}j^{r-1}\biggr)=\\


%


=r\biggl(\,\sum^n_{j=1}j^{r-1}\biggr)\varphi_n+


r\sum^\infty_{\nu=n+1}\Delta\varphi_{\nu-1}\sum^\nu_{j=n+1}


j^{r-1}\le rn^r\varphi_n+r\sum^\infty_{j=n+1}j^{r-1}\varphi_{j-1}.


\end{multline*}


Соединяя (23)--(25), имеем для $l>r\ge 1$


$$


(l+o(1)-r)\sum^\infty_{\nu=n+1}\nu^{r-1}\varphi_{\nu-1}


-rn^r\varphi_n=O(E_n(f^{(r)}_3)).


$$


Отсюда ввиду $cf_3\in H_l(\omega)$ по функциям $f\in H_l(\omega)$


получаем соотношение


\begin{equation}


\sum^\infty_{\nu=n+1}\nu^{r-1}\varphi_{\nu-1}=O(\sup E_n(f^{(r)})+


n^r\varphi_n),\quad \varphi_\nu=\omega^{**}_l(\nu^{-1}).


\end{equation}


Остается доказать, что


\begin{equation}


n^r\varphi_n=O(\sup E_n(f^{(r)})),\quad f\in H_l(\omega).


\end{equation}


Но это ясно из рассмотрения функции


$f_0(x)=\varphi_n\cos(n+1)x$, для которой очевидно имеем


$\omega_l(f_0,m^{-1})=O(\varphi_m)$ и


$E_n(f^{(r)}_0)=n^r\varphi_n$. Сопоставляя (26), (27), где


$\varphi_n=\omega^{**}_l(n^{-1})$, из (20)


получим теорему 5 при $r$ нечетном.





Пусть теперь $r$ четное, тогда по предположению теоремы 5 данное $l$


должно


быть нечетным. В этом случае по заданной мажоранте


$\omega(\delta)$ с $\omega^{**}_l(\delta)\not\equiv 0$ рассмотрим


функцию $f_0$, которая получается из (3) доказательства теоремы 1 при


$r=0$ и


с заменой там $F_\nu$ на $\omega^{**}_l(\nu^{-1})$. Так как $l$ нечетное,


то по лемме 2 (\cite{Ge1},


с.\,71) имеем $\omega_l(f_0,n^{-1})=O(\omega^{**}_l(n^{-1}))$. С другой


стороны, производная $f^{(r)}_0$ в силу четности


$r$ будет $\cos$-рядом, для которого имеет место соотношение (6). Значит,


поскольку $F_\nu=\omega^{**}_l(\nu^{-1})$,


$$


\sum^\infty_{\nu=n+1}\nu^{r-1}\omega^{**}_l(\nu^{-1})


=O(E_n(f^{(r)}_0))


$$


(см. текст, сопровождающий (6) в доказательстве теоремы 1).


Теорема 5 доказана.





Отправляясь теперь от неравенства (1.7) из


(\cite{Bar}, c.\,488) и привлекая соответствующим образом средства, уже


использованные при доказательстве оценок снизу в теореме 5, получим ее


аналог для сопряженных функций в следующем виде.





\begin{thm}


Пусть заданы натуральные $l>r$ и мажоранта $\omega(\delta)$ с


$\omega^{**}_l(\delta)\not\equiv 0$, тогда


кроме случая четного $l$ при $r$ нечетном имеем


$$


\sup E_n(\Tilde f^{(r)})\sim \sum^\infty_{\nu=n+1}\nu^{r-1}


\omega^{**}_l(\nu^{-1}),\quad f\in H_l(\omega),


$$


в предположении сходимости указанного здесь ряда.


\end{thm}





\begin{thm}


Пусть заданы натуральные $k$, $l$, $r$ и мажоранта


$\omega(\delta)$ с $\omega^{**}_l(\delta)\not\equiv 0$, причем


$k+r>l>r$,


тогда кроме случая, когда $r$, $l$ оба четные, имеем


\begin{equation}


\sup\omega_k(f^{(r)},n^{-1})\sim \sum^\infty_{\nu=n+1}


\nu^{r-1}\omega^{**}_l(\nu^{-1}),\quad f\in H_l(\omega),


\end{equation}


при условии сходимости ряда в {\em (28)}.


\end{thm}





\begin{pf}


Сопоставляя на этот раз обобщенное неравенство Джексона


$E_\nu(f)=O(\omega_l(f,\nu^{-1}))$ с соотношением


(9), получим


\begin{equation}


\omega_k(f^{(r)},n^{-1})=O\biggl(n^{-k}


\sum^n_{\nu=1}\nu^{k+r-1}\omega_l(f,\nu^{-1})+\sum^\infty_{\nu=n+1}


\nu^{r-1}\omega_l(f,\nu^{-1})\biggr).


\end{equation}


Пусть $f\in H_l(\omega)$, тогда


$\omega_l(f,\nu^{-1})=O(\omega^{**}_l(\nu^{-1}))$ \cite{St}, поэтому для


таких $f$ имеем


\begin{equation}


\omega_k(f^{(r)},n^{-1})=O\biggl(n^{-k}\sum^n_{\nu=1}


\nu^{k+r-1}\omega^{**}_l(\nu^{-1})+\sum^\infty_{\nu=n+1}


\nu^{r-1}\omega^{**}_l(\nu^{-1})\biggr).


\end{equation}


Так как по (18) последовательность


$\nu^l\omega^{**}_l(\nu^{-1})$ не убывает при возрастании


$\nu\uparrow\infty$, то


для первого слагаемого в (30) при $k+r>l$ будем иметь последовательно


$$


n^{-k}\sum^n_{\nu=1}\nu^{k+r-l-1}\nu^l\omega^{**}_l(\nu^{-1})\le \\


n^{-k}n^l\omega^{**}_l(n^{-1})\sum^n_{\nu=1}


\nu^{k+r-l-1}=O(n^r\omega^{**}_l(n^{-1})).


$$


Нетрудно показать далее, что величина


$n^r\omega^{**}_l(n^{-1})$ есть $O$ от второго слагаемого в


(30)\footnote{Аналогично устраняется первое слагаемое и в (29) для


$k+r>l$.}.


Поэтому при $k+r>l$


\begin{equation}


\sup\omega_k(f^{(r)},n^{-1})=O\biggl(\,


\sum^\infty_{\nu=n+1}\nu^{r-1}\omega^{**}_l(\nu^{-1})\biggr),\quad


f\in H_l(\omega),


\end{equation}


в предположении сходимости этого ряда и независимо от четности


$r$, $l$.





Если же исключить случай, когда $r$, $l$ оба четные, то по теореме 5 и


благодаря обобщенному неравенству Джексона имеем обратное к (31)


соотношение.


\end{pf}





Отправляясь на этот раз от неравенства (1.19) из (\cite{Bar}, c.\,492),


имеем


\begin{equation}


\omega_k(\Tilde f^{(r)},n^{-1})=O\biggl(n^{-k}\sum^n_{\nu=1}


\nu^{k+r-1}\omega_l(f,\nu^{-1})+\sum^\infty_{\nu=n+1}


\nu^{r-1}\omega_l(f,\nu^{-1})\biggr),


\end{equation}


где $l>r$. Как и выше (см. сноску), преобразуем (32) при


$k+r>l$ к виду


\begin{equation}


\omega_k(\Tilde f^{(r)},n^{-1})=O\biggl(\,\sum^\infty_{\nu=n+1}


\nu^{r-1}\omega_l(f,\nu^{-1})\biggr).


\end{equation}


Отсюда и при помощи приведенной выше теоремы 6 получается следующая


теорема для сопряженных функций.





\begin{thm}


Пусть заданы натуральные $k$, $l$, $r$ и мажоранта


$\omega(\delta)$ с $\omega^{**}_l(\delta)\not\equiv 0$, причем


$k+r>l>r$.


Тогда, кроме случая, когда $l$ четное, а $r$ нечетное, имеем


$$


\sup\omega_k(\Tilde f^{(r)},n^{-1})\sim


\sum^\infty_{\nu=n+1}\nu^{r-1}\omega^{**}_l(\nu^{-1}),\quad


f\in H_l(\omega),


$$


в предположении сходимости указанного здесь ряда.


\end{thm}





Условия на $k$, $l$, $r$ в нижеследующих теоремах таковы, что


доказательства не потребуют привлечения новых средств, кроме уже


использованных выше. Поэтому уместно ограничиться здесь одними лишь


формулировками теорем.





\begin{thm}


Пусть заданы натуральные $k$, $l$, $r$ и мажоранта


$\omega(\delta)$, причем $k+r\le l$ и


сходится ряд в {\em (34)}, тогда


\begin{equation}


\sup\omega_k(f^{(r)},n^{-1})\sim n^{-k}


\sum^n_{\nu=1}\nu^{k+r-1}\omega^{**}_l(\nu^{-1})+


\sum^\infty_{\nu=n+1}\nu^{r-1}\omega^{**}_l(\nu^{-1}),\quad


f\in H_l(\omega),


\end{equation}


в случаях{\em :} {\em (a)} оба числа $l$, $k$ четные при $r$


нечетном{\em ;}


{\em (b)} $l$ нечетное, а $k+r$ четное.


\end{thm}





\begin{thm}


Пусть заданы натуральные $k$, $l$, $r$ и мажоранта


$\omega(\delta)$, причем $k+r< l$ и


сходится ряд в {\em (35)}, тогда


\begin{equation}


\sup\omega_k(\Tilde f^{(r)},n^{-1})\sim n^{-k}


\sum^n_{\nu=1}\nu^{k+r-1}\omega^{**}_l(\nu^{-1})+


\sum^\infty_{\nu=n+1}\nu^{r-1}\omega^{**}_l(\nu^{-1}), \quad


f\in H_l(\omega),


\end{equation}


в случаях{\em :} {\em (a)} числа $l$, $r$, $k$ все четные{\em ;}


{\em (b)} оба числа $l$, $k+r$ нечетные.


\end{thm}





\begin{notenonum}


Задача о порядке рассматривавшихся здесь величин имеет отчасти смысл и


при $r=0$. Соответствующие случаи подробно изучались автором в работе


\cite{Ge2}, где излагается также история данного вопроса.


\end{notenonum}








\begin{thebibliography}{9}





\bibitem{Bar}


Бари Н.К., Стечкин С.Б. {\em Наилучшие приближения и дифференциальные


свойства двух сопряженных функций} // Тр. Моск. матем. о-ва. -- 1956. --


Т.\,5. -- С.\,483--522.


\bibitem{Cen}


Чень Тянь-Пинь. {\em О тригонометрических рядах со знакоопределенными


коэффициентами} // Acta scient. natur. Univ. Fudan. -- 1966. -- V.\,11. --


\No\,1. -- P.\,1--6.


\bibitem{Ge1}


Гейт В.Э. {\em Теоремы вложения для некоторых классов периодических


непрерывных функций} // Изв. вузов. Математика. -- 1972. -- \No\,4. --


С.\,67--77.


\bibitem{St}


Стечкин С.Б. {\em Об абсолютной сходимости рядов Фурье} // Изв. АН СССР.


Сер. матем. -- 1955. -- Т.\,19. -- \No\,4. -- С.\,221--246.


\bibitem{Ge2}


Гейт В.Э. {\em Структурные и конструктивные свойства функции и ее


сопряженной}\,: Дис.$\;\dots$ канд. физ.-матем. наук. -- Свердловск, 1970.


-- 148 с.


\bibitem{Ge3}


Гейт В.Э. {\em Об условиях вложения классов $H^\omega_{k,R}$ и


$\widetilde H^\omega_{k,R}$} // Матем. заметки. -- 1973. -- Т.\,13. --


\No\,2. -- C.\,169--178.





\end{thebibliography}





\vskip 2cm





\post{\it Челябинский государственный университет}{\it Поступила}


\post{}{24.01.1995}





\end{document}


